\chapter{The DAGonStar Computational Model and Architecture}
\label{ch:architecture}

This chapter develops, in precise and implementation-faithful terms, the
internal architecture introduced informally in \Cref{ch:introduction}. The
DAGonStar documentation states a strict principle that this chapter respects
throughout: \textit{``documentation must remain consistent with the software
in this repository''} --- this is a descriptive account of the current
implementation, not an aspirational design document.

\section{Package Layout}

The engine's source tree is organized as follows:

\begin{lstlisting}[language=bash, caption={\texttt{dagon/} package layout.}]
dagon/
  __init__.py                 Workflow, Status, DataMover, Stager
  task.py                     TaskType, DagonTask, Task base class
  batch.py                    Batch, RemoteBatch, Slurm, RemoteSlurm
  remote.py                   RemoteTask, CloudTask
  docker_task.py              DockerTask, DockerRemoteTask
  checkpoint.py               Checkpoint tasks
  config.py                   INI configuration reader
  dag_tps.py                  Meta-workflow coordination
  api/                        HTTP client and workflow service
  cloud/                      Apache Libcloud helpers
  communication/               SSH, SCP, Globus, SKYCDS helpers
  ftp_publisher/               FTP publishing helper
  peer2peer/                   Experimental peer-to-peer service functions
\end{lstlisting}

\begin{figure}[h]
\centering
\begin{tikzpicture}[
  layer/.style={rectangle, draw=dagonblue, very thick, rounded corners,
    minimum width=120mm, minimum height=10mm, font=\sffamily\small},
  fill1/.style={fill=blue!4}, fill2/.style={fill=blue!9},
  fill3/.style={fill=blue!14}, fill4/.style={fill=blue!20},
]
  \node[layer, fill1] (a) at (0,0)  {\textbf{Authoring layer} --- \texttt{Workflow}, \texttt{DagonTask}, \texttt{DAG\_TPS}};
  \node[layer, fill2] (b) at (0,-1.3) {\textbf{Dependency \& scheduling layer} --- \texttt{Task} (threaded), \texttt{Status}, \texttt{workflow://} resolution};
  \node[layer, fill3] (c) at (0,-2.6) {\textbf{Staging layer} --- \texttt{Stager}, \texttt{DataMover}, \texttt{StagerMover}};
  \node[layer, fill4] (d) at (0,-3.9) {\textbf{Execution backends} --- \texttt{Batch}, \texttt{Slurm}, \texttt{RemoteTask}, \texttt{CloudTask}, \texttt{DockerTask}, \texttt{LLMTask}, \texttt{NativeTask}, \texttt{WebTask}};
  \node[layer, fill1] (e) at (0,-5.2) {\textbf{Optional services} --- \texttt{dagon.api} (Flask service), \texttt{cloud/}, \texttt{communication/} (SSH/SCP/Globus/SKYCDS), \texttt{ftp\_publisher/}};
\end{tikzpicture}
\caption{Layered view of the DAGonStar architecture. Each layer depends only
on the layer(s) below it; optional services are loaded lazily so that core
local-batch workflows never require Docker, cloud, or Globus dependencies to
be installed.}
\label{fig:layered-architecture}
\end{figure}

\section{Main Abstractions}

\subsection{\texttt{Workflow}}

\texttt{dagon.Workflow} owns a list of tasks and coordinates dependency
construction, execution, checkpoint state, configuration, and optional
service registration. Its most important methods are:

\begin{itemize}[leftmargin=1.4em]
  \item \texttt{add\_task(task)} --- register a task with the workflow.
  \item \texttt{make\_dependencies()} --- parse \texttt{workflow://}
    references and validate the resulting graph.
  \item \texttt{run(resume\_checkpoint\_file=None)} --- execute the workflow
    synchronously, in the calling thread.
  \item \texttt{launch(resume\_checkpoint\_file=None)} and
    \texttt{wait(timeout=None)} --- execute asynchronously, in a background
    thread (\Cref{ch:meta-async}).
  \item \texttt{as\_json()} --- serialize the workflow's structure.
  \item \texttt{find\_task\_by\_name(workflow\_name, task\_name)} --- locate a
    task, possibly across workflow boundaries.
  \item \texttt{Validate\_WF()} --- reject cyclic graphs (\Cref{sec:dag-theory}).
  \item \texttt{set\_data\_mover(data\_mover)} /
    \texttt{set\_stager\_mover(stager\_mover)} --- configure how files move
    between tasks.
\end{itemize}

\subsection{\texttt{DagonTask}: the Task Factory}

\texttt{DagonTask} is a \emph{factory}: it receives a \texttt{TaskType} value
and returns an instance of the concrete task class mapped to that type.

\begin{lstlisting}[caption={\texttt{DagonTask} as a factory over concrete task classes.}]
task = DagonTask(TaskType.BATCH, "A", "echo A")
\end{lstlisting}

This single line hides the dispatch logic: depending on \texttt{TaskType} and
on auxiliary keyword arguments (such as \texttt{ip} for remote execution),
\texttt{DagonTask} returns the appropriate subclass without the workflow
author needing to import or instantiate it directly.

\subsection{\texttt{Task}: the Common Base}

\texttt{Task} extends \texttt{threading.Thread} and provides, uniformly
across every execution backend:

\begin{itemize}[leftmargin=1.4em]
  \item task status (see \Cref{sec:status-model});
  \item predecessor (\texttt{prevs}) and successor (\texttt{nexts}) lists;
  \item a back-reference to the owning \texttt{Workflow};
  \item scratch-directory creation;
  \item \texttt{workflow://} parsing (\Cref{ch:schema});
  \item launcher-script generation;
  \item reference counting and garbage-collection hooks;
  \item checkpoint interaction (\Cref{ch:lessons-05-07}, Appendix~D).
\end{itemize}

Concrete task classes override \texttt{on\_execute()} and related
environment-specific hooks to specialize \emph{how} a command is actually
run, while inheriting all dependency, staging, and checkpoint behavior from
\texttt{Task}. This is the architectural reason a heterogeneous DAG --- mixing
a local shell command, a Slurm submission, and an LLM completion --- behaves
uniformly with respect to dependencies and data movement: heterogeneity lives
only in \texttt{on\_execute()}.

\section{Execution Lifecycle}
\label{sec:execution-lifecycle}

A typical dataflow run follows a fifteen-step sequence, traced here exactly
as the architecture documentation specifies it:

\begin{enumerate}[leftmargin=1.6em]
  \item The user creates a \texttt{Workflow}.
  \item The user creates tasks through \texttt{DagonTask}.
  \item The user adds tasks with \texttt{workflow.add\_task()}.
  \item The user calls \texttt{workflow.make\_dependencies()}.
  \item Each task parses its command for \texttt{workflow://} references.
  \item Dependencies are recorded in \texttt{prevs} and \texttt{nexts}.
  \item \texttt{Workflow.Validate\_WF()} rejects cycles.
  \item The user calls \texttt{workflow.run()}.
  \item Tasks are started as Python threads.
  \item Each task waits for predecessor completion.
  \item Each task creates or reuses a scratch directory.
  \item Input data is staged.
  \item The task launcher script is executed.
  \item Status and checkpoint metadata are updated.
  \item Dependent tasks proceed.
\end{enumerate}

\begin{figure}[h]
\centering
\begin{tikzpicture}[
  node distance=7mm,
  lstep/.style={rectangle, draw=dagonblue, rounded corners, fill=blue!4,
    text width=46mm, align=center, font=\sffamily\scriptsize,
    minimum height=8mm},
  arr/.style={-{Stealth[length=2.2mm]}, thick, dagonblue}
]
  \node[lstep] (s1) {1--3. Author workflow \& tasks, \texttt{add\_task()}};
  \node[lstep, below=of s1] (s2) {4--7. \texttt{make\_dependencies()}: parse \texttt{workflow://}, build \texttt{prevs}/\texttt{nexts}, \texttt{Validate\_WF()}};
  \node[lstep, below=of s2] (s3) {8--9. \texttt{run()} starts tasks as threads};
  \node[lstep, below=of s3] (s4) {10. Each task waits on predecessors};
  \node[lstep, below=of s4] (s5) {11--12. Scratch dir created; inputs staged};
  \node[lstep, below=of s5] (s6) {13. Launcher script executed};
  \node[lstep, below=of s6] (s7) {14--15. Status/checkpoint updated; dependents released};
  \draw[arr] (s1) -- (s2);
  \draw[arr] (s2) -- (s3);
  \draw[arr] (s3) -- (s4);
  \draw[arr] (s4) -- (s5);
  \draw[arr] (s5) -- (s6);
  \draw[arr] (s6) -- (s7);
  \draw[arr] (s7.east) .. controls +(2.4,0) and +(2.4,0) .. node[right=1mm,
    font=\sffamily\scriptsize, text width=18mm] {next ready task} (s4.east);
\end{tikzpicture}
\caption{The fifteen-step execution lifecycle, grouped into seven logical
phases. The loop at the right indicates that step~15 (releasing dependents)
re-enters the per-task lifecycle for every task whose predecessors have just
completed.}
\label{fig:execution-lifecycle}
\end{figure}

\section{The Status Model}
\label{sec:status-model}

\texttt{dagon.Status} defines five task states:

\begin{center}
\begin{tabular}{@{}ll@{}}
\toprule
\textbf{Status} & \textbf{Meaning} \\
\midrule
\texttt{READY}    & Eligible to run; no unmet dependency. \\
\texttt{WAITING}  & Blocked on at least one predecessor. \\
\texttt{RUNNING}  & Launcher script currently executing. \\
\texttt{FINISHED} & Completed with a successful exit code. \\
\texttt{FAILED}   & Completed with a non-zero exit code or raised exception. \\
\bottomrule
\end{tabular}
\end{center}

Status changes are local to the running process and, when the optional
DAGon service is enabled, are additionally propagated through the service
API (\Cref{sec:service-architecture}).

\section{The Dependency Model}

Every task carries three pieces of dependency bookkeeping:

\begin{itemize}[leftmargin=1.4em]
  \item \texttt{prevs}: the tasks that must finish before this task may run;
  \item \texttt{nexts}: the tasks that depend on this task;
  \item \texttt{reference\_count}: the number of downstream data references,
    used by garbage-collection hooks for scratch-directory cleanup.
\end{itemize}

Explicit dependencies are created with \texttt{consumer.add\_dependency\_to(producer)};
implicit dependencies are created automatically by parsing
\texttt{workflow:///producer/path/to/file} tokens in a task's command, as
introduced in \Cref{sec:task-vs-data} and developed fully in
\Cref{ch:schema}.

\section{Data Staging}
\label{sec:data-staging-architecture}

\texttt{Stager.stage\_in()} selects a data-movement strategy for a given
dependency edge. The implemented \texttt{DataMover} enumeration is:

\begin{center}
\begin{tabular}{@{}ll@{}}
\toprule
\textbf{\texttt{DataMover}} & \textbf{Mechanism} \\
\midrule
\texttt{DONTMOVE} & No staging is performed. \\
\texttt{LINK}     & Symbolic link to the producer's file. \\
\texttt{COPY}     & Independent file copy. \\
\texttt{SCP}      & Secure copy over SSH. \\
\texttt{HTTP} / \texttt{HTTPS} & Movement over plain or secured HTTP. \\
\texttt{FTP} / \texttt{SFTP}   & Movement over (S)FTP. \\
\texttt{GRIDFTP}  & GridFTP-based movement (Globus). \\
\texttt{SKYCDS}   & SKYCDS-backed catalog/data movement. \\
\bottomrule
\end{tabular}
\end{center}

The currently most exercised local paths are \texttt{COPY} and \texttt{LINK};
remote and specialized movers require external services and site-specific
configuration. A complementary enumeration, \texttt{StagerMover}, controls
the \emph{staging mode}: \texttt{NORMAL}, \texttt{PARALLEL}, or
\texttt{SLURM}.

\section{Execution Backends}

\subsection{Batch}
\texttt{Batch} executes a generated launcher script locally through
\texttt{subprocess.Popen}. This is the default, dependency-free execution
path used by nearly every example in this book.

\subsection{Remote Batch}
When a batch task is constructed with \texttt{ip=...}, the factory
\texttt{Batch.\_\_new\_\_} transparently returns a \texttt{RemoteBatch}
instance instead, which executes the same generated launcher through an
\texttt{SSHManager}-mediated connection.

\subsection{Slurm}
\texttt{Slurm} generates an \texttt{sbatch} command from optional parameters
such as partition, task count, memory, wall time, nodes, and tasks per node
--- the natural interface for compiled HPC models (\Cref{ch:backends}).

\subsection{Docker}
\texttt{DockerTask} creates or reuses a Docker container through the Docker
Python SDK; the task command is wrapped inside \texttt{docker exec}.

\subsection{Cloud}
\texttt{CloudTask} extends remote execution with Apache Libcloud-backed
instance management, allowing a task to provision or reuse a cloud virtual
machine as its execution target.

\subsection{Checkpoint}
Checkpoint tasks preserve task working directories and checkpoint metadata.
The base \texttt{Workflow} can resume from a checkpoint file independently of
any explicit \texttt{CHECKPOINT} task (Appendix~D).

\section{Service Architecture}
\label{sec:service-architecture}

The optional \texttt{dagon.api} package contains two components:

\begin{itemize}[leftmargin=1.4em]
  \item \texttt{API}: a client used by workflows to register and update
    remote service state.
  \item \texttt{WorkflowServer}: a Flask-based server wrapper exposing
    workflow operations over HTTP.
\end{itemize}

The default, documented configuration disables this service entirely:

\begin{lstlisting}[style=ini]
[dagon_service]
use = False
\end{lstlisting}

When disabled --- the recommended mode for introductory examples, unit tests,
and the entirety of this textbook's local exercises --- DAGonStar runs as a
purely local, in-process workflow engine with no external service
dependency.

\section{Known Architectural Constraints}
\label{sec:known-constraints}

In the interest of scientific and engineering honesty, this section
reproduces, verbatim in substance, the constraints the project itself
documents as open or partially mitigated:

\begin{cautionbox}
\begin{itemize}[leftmargin=1.4em]
  \item Tasks are Python \emph{threads}, not isolated operating-system
    processes.
  \item Some command strings are still assembled by string concatenation;
    shell-command safety and path quoting remain an active improvement
    priority, although new staging helpers already quote generated staging
    paths.
  \item Optional integrations (Docker, cloud, Globus, API) are exposed
    through package \emph{extras} and use lazy imports, so that core
    local-batch workflows do not require those dependencies to be installed.
  \item The repository has an initial unit-test and CI baseline, but
    comprehensive workflow, checkpoint, staging, and integration coverage is
    still being expanded.
  \item Public APIs have broad type annotations across workflow, task,
    configuration, staging, task-subclass, and API-client code; stricter
    static checking remains a maintainability goal.
  \item Some examples depend on external systems and cannot be verified in a
    generic continuous-integration environment.
\end{itemize}
\end{cautionbox}

\section{The Factory Pattern in Depth}
\label{sec:factory-pattern-depth}

\texttt{DagonTask} deserves a closer reading than \Cref{sec:computational-model}
gave it, because the factory pattern it implements is the architectural
device that lets Part~II of this book introduce four further execution
backends (Docker, remote/Slurm, LLM, Web, Native) without ever changing the
authoring surface a workflow writer interacts with.

In classical object-oriented design terms, a \emph{factory} is a function or
method that returns an instance of one of several related classes, selected
by a parameter, without the caller needing to know --- or import --- the
concrete class being instantiated. \texttt{DagonTask(TaskType.BATCH, ...)}
and \texttt{DagonTask(TaskType.SLURM, ...)} return instances of entirely
different classes (\texttt{Batch} and \texttt{Slurm} respectively, per the
\texttt{tasks\_types} mapping reproduced in Appendix~A), yet both are
constructed through an identical calling convention, and both are usable
interchangeably as nodes in a \texttt{Workflow}'s task list. This is what
makes the heterogeneous DAGs explored from \Cref{ch:backends} onward
possible at all: a single \texttt{workflow.add\_task()} call accepts any
concrete task type with no special-casing required at the workflow level.

The remote-variant dispatch described in \Cref{ch:architecture} --- a
\texttt{BATCH} task constructed with \texttt{ip=...} silently becoming a
\texttt{RemoteBatch} --- is a second, nested instance of the same pattern,
implemented one level lower, inside \texttt{Batch.\_\_new\_\_} rather than
inside \texttt{DagonTask} itself. The practical lesson for a workflow author
is that \emph{keyword arguments, not just the \texttt{TaskType} value, are
part of the dispatch key}: two superficially similar \texttt{DagonTask}
calls that differ only in the presence of an \texttt{ip} keyword can
construct objects of entirely different concrete classes, executing in
entirely different environments, while remaining structurally identical from
the perspective of the dependency graph.

\section{Threads, Not Processes: Consequences and Caveats}
\label{sec:threads-not-processes}

\Cref{sec:known-constraints} flagged, among the documented architectural
constraints, that DAGonStar tasks are Python \texttt{threading.Thread}
objects rather than isolated operating-system processes. This single design
decision has consequences pervasive enough to merit a dedicated discussion,
because it shapes what a workflow author may safely assume about isolation,
concurrency, and resource accounting.

In CPython, the standard implementation most DAGonStar deployments will run
on, the Global Interpreter Lock (GIL) serializes execution of Python
bytecode across threads within a single process: at any instant, only one
thread is actually executing Python code, even on a multi-core machine. This
does \emph{not}, however, prevent DAGonStar's concurrent task execution from
delivering real parallelism for the cases that matter most in scientific
computing, for a structural reason: a \texttt{Batch} task's actual
computational work happens inside a \emph{child process} spawned via
\texttt{subprocess.Popen}, entirely outside the GIL's purview. The
controlling Python thread for that task spends almost all of its lifetime
blocked waiting for the child process to exit --- a blocking I/O-bound wait,
which CPython threads handle well, releasing the GIL for other threads while
blocked. The same reasoning extends to \texttt{RemoteTask} (blocked on a
network socket during SSH execution), \texttt{DockerTask} (blocked on the
Docker daemon's API), \texttt{WebTask} (blocked on an HTTP response), and
\texttt{LLMTask} (blocked on the same). Only \texttt{NativeTask}, whose
computation runs as actual in-process Python code (\Cref{ch:native-tasks}),
is potentially GIL-bound, and therefore is the one task type whose
\emph{computational}, rather than purely I/O, concurrency benefits most
directly from the engine's \texttt{executor="slurm"} escape hatch to
out-of-process execution.

\begin{cautionbox}
Thread-based isolation is weaker than process-based isolation in ways that
matter for defensive workflow design. Threads share the controlling Python
process's memory space, open file descriptors, environment variables, and
current working directory at the moment of thread creation. A bug in one
task's Python-level pre/post-processing code --- as opposed to the external
program or subprocess it launches --- could in principle corrupt shared
interpreter state in a way a process-isolated design would prevent
entirely. This is part of why the developer guide (Appendix~B) flags
defensive shell quoting and careful handling of paths and user-controlled
values as an active, ongoing engineering priority rather than a solved
problem.
\end{cautionbox}

\section{Comparing the Layered Architecture to Alternative Designs}

It is instructive to contrast the four-layer architecture of
\Cref{fig:layered-architecture} with two alternative architectural strategies
a workflow engine could have adopted, because the contrast clarifies what
DAGonStar is optimizing for.

A \textbf{monolithic engine} would fold dependency resolution, staging, and
execution into a single class hierarchy, with each concrete task type
re-implementing its own version of dependency tracking and data movement.
This would reduce the indirection a reader must follow to understand any one
task type in isolation, but it would multiply the surface area for bugs in
dependency resolution --- a defect in cycle detection, for instance, would
need to be fixed and tested independently in every concrete task class,
rather than once in the shared \texttt{Task} base class examined in
\Cref{sec:execution-lifecycle}.

A \textbf{microservice-oriented engine}, at the opposite extreme, would
externalize scheduling and staging into separate, independently deployed
services communicating over a network protocol, with the Python library
reduced to a thin client. This would more cleanly separate concerns and
would scale naturally to very large workflows or multi-tenant deployments,
but at the cost of introducing a mandatory network dependency and
operational complexity (service discovery, authentication, availability)
for even the simplest single-machine workflow --- exactly the cost the
optional, disable-by-default DAGon service architecture
(\Cref{sec:service-architecture}) is designed to avoid imposing on the
common case.

\begin{practicebox}
DAGonStar's layered-but-in-process design (\Cref{fig:layered-architecture})
occupies a deliberate middle position: dependency resolution and staging are
shared, reusable layers (avoiding the monolithic engine's duplication), but
they execute in-process by default (avoiding the microservice engine's
mandatory operational overhead), with the optional \texttt{dagon.api}
package available as an escape hatch when a deployment genuinely needs
multi-process or multi-host coordination beyond what SSH-based remote
execution already provides.
\end{practicebox}

\section{Reading the Architecture Document as a Living Artifact}

A final, methodological point belongs in this chapter rather than in the
developer guide (Appendix~B), because it concerns how a \emph{reader},
not only a contributor, should treat architectural documentation.
\Cref{ch:architecture} opened by quoting the project's own principle that
documentation must remain consistent with the current implementation, and
\Cref{sec:known-constraints} reproduced, deliberately without softening, the
project's own list of open architectural debt. This is not incidental
even-handedness; it reflects a specific, defensible position about what
makes technical documentation trustworthy for a scientific audience.

A textbook --- or any documentation --- that described only an engine's
intended, idealized behavior, omitting known gaps such as incomplete shell
quoting or partial test coverage, would actively mislead a researcher
relying on it to assess whether a given workflow's results are
trustworthy. The constraints listed in \Cref{sec:known-constraints} are not
flaws unique to DAGonStar; comparable caveats (partial coverage of edge
cases, ongoing hardening of command construction, uneven test coverage
across optional integrations) are common to research software more broadly,
precisely because research software is typically developed under different
incentives than commercial infrastructure software. Treating these caveats
as part of the architecture, rather than as a separate errata list, is this
book's attempt to model the kind of honest engineering communication the
developer guide (Appendix~B) explicitly asks contributors to maintain.

\section{Chapter Summary}

The architecture is a thin, layered stack: an authoring layer
(\texttt{Workflow}, \texttt{DagonTask}), a dependency and scheduling layer
built on threaded \texttt{Task} objects, a staging layer that abstracts data
movement behind \texttt{DataMover}/\texttt{StagerMover}, a family of
execution backends that share this common infrastructure while specializing
only \texttt{on\_execute()}, and an optional service layer for distributed
registration and status propagation. This chapter additionally examined the
factory pattern underlying \texttt{DagonTask} and its nested remote-variant
dispatch, reasoned about the consequences of thread-based rather than
process-based task isolation, compared the chosen layered-in-process design
against monolithic and microservice alternatives, and argued for treating
architectural documentation --- including its documented limitations --- as
part of the scientific record a researcher relies on. The next chapter puts
this architecture to work by walking through installation and a complete
first workflow.

\begin{exercise}
Explain, in terms of \Cref{sec:threads-not-processes}, why a
\texttt{NativeTask} performing a computationally intensive, pure-Python
operation (for example, a nested loop over a large array with no NumPy
vectorization) might \emph{not} benefit from running concurrently with
several sibling \texttt{Batch} tasks in the same workflow, even though the
engine schedules all of them as threads. What change to the task's
implementation, or to its \texttt{executor} argument, would most directly
address this?
\end{exercise}
